% p027_d 的电脑重画（P外 随 I、U、R 变化的三幅图像：P外-I、P外-U 为过原点的抛物线（峰值 E²/4r），P外-R 先升后降、峰在 R=r）
\begin{tikzpicture}[line width=0.7pt, >={Stealth[length=2.2mm]}]
  % (1) P外 - I
  \draw[->] (0,0) -- (3.0,0) node[right] {$I$};
  \draw[->] (0,0) -- (0,2.6) node[above left] {$P_{\text{外}}$};
  \draw[domain=0:2.3, samples=40, smooth] plot (\x, {7.2*(\x/2.3)*(1-\x/2.3)});
  \draw[dashed] (0,1.8) -- (1.15,1.8) -- (1.15,0);
  \node[left] at (0,1.8) {$\dfrac{E^2}{4r}$};
  \node[below] at (1.15,0) {$\dfrac{I_0}{2}$};
  \node[below] at (2.3,0) {$I_0$};
  \node[anchor=north west] at (2.1,-0.95) {短路电流};
  % (2) P外 - U
  \begin{scope}[shift={(4.6,0)}]
    \draw[->] (0,0) -- (3.0,0) node[right] {$U$};
    \draw[->] (0,0) -- (0,2.6) node[above left] {$P_{\text{外}}$};
    \draw[domain=0:2.0, samples=40, smooth] plot (\x, {7.2*(\x/2.0)*(1-\x/2.0)});
    \draw[dashed] (0,1.8) -- (1.0,1.8) -- (1.0,0);
    \node[left] at (0,1.8) {$\dfrac{E^2}{4r}$};
    \node[below] at (1.0,0) {$\dfrac{E}{2}$};
    \node[below] at (2.0,0) {$E$};
  \end{scope}
  % (3) P外 - R
  \begin{scope}[shift={(9.2,0)}]
    \draw[->] (0,0) -- (3.3,0) node[right] {$R$};
    \draw[->] (0,0) -- (0,2.6) node[above left] {$P_{\text{外}}$};
    \draw (0,0) .. controls (0.55,1.1) and (0.85,1.8) .. (1.3,1.8)
          .. controls (1.7,1.8) and (1.95,1.3) .. (2.3,0.65)
          .. controls (2.5,0.28) and (2.75,0.25) .. (3.1,0.25);
    \draw[dashed] (0,1.8) -- (1.3,1.8) -- (1.3,0);
    \node[left] at (0,1.8) {$\dfrac{E^2}{4r}$};
    \node[below] at (1.3,0) {$R=r$};
  \end{scope}
\end{tikzpicture}
